% Author's solution, images/HW04/IMG_9194.HEIC--IMG_9197.HEIC.
% The derivative's cubic coefficient is 3a, as in IMG_9196's final matrix;
% IMG_9194--IMG_9195's 2a conflicts with that final result.
% IMG_9195's edited conclusion says D is surjective.
\begin{proof}[Solution]
\begin{enumerate}[(a)][2]
  \item Write
    \[
      p(x)=ax^3+bx^2+cx+d.
    \]
    Then \(D(p)=p'=3ax^2+2bx+c\), which is zero exactly when
    \(a=b=c=0\). The constant coefficient \(d\) is unrestricted.
    Therefore
    \[
      \ker D=\operatorname{span}\{1\},
      \qquad
      \operatorname{nullity}D=1.
    \]

    Every derivative belongs to \(\mathcal P_2(\mathbb{R})\). Conversely, given
    \(q(x)=ux^2+vx+w\), choose
    \[
      p(x)=\frac{u}{3}x^3+\frac{v}{2}x^2+wx.
    \]
    Then \(D(p)=q\). Hence
    \[
      \operatorname{im}D=\mathcal P_2(\mathbb{R})
        =\operatorname{span}\{1,x,x^2\},
      \qquad
      \operatorname{rank}D=3.
    \]

  \item The map \(D\) is not injective because its nullity is \(1\).
    It is surjective because its image is all of \(\mathcal P_2(\mathbb{R})\).

  \item Apply \(D\) to the given basis entries. We obtain
    \[
      D(E)=[\,D(1)\ \ D(x)\ \ D(x^2)\ \ D(x^3)\,]
        =[\,0\ \ 1\ \ 2x\ \ 3x^2\,].
    \]
    Expressing these entries in \(F=[\,1\ \ x\ \ x^2\,]\) gives
    \[
      M=
      \begin{pmatrix}
        0&1&0&0\\
        0&0&2&0\\
        0&0&0&3
      \end{pmatrix}.
    \]
    Thus \(D(E)=FM\).

  \item We want the nonzero basis images to be \(1,x,x^2\), respectively.
    Taking their antiderivatives gives
    \[
      E'=\left[\,1\ \ x\ \ \frac{x^2}{2}\ \ \frac{x^3}{3}\,\right].
    \]
    Its first vector is a basis of \(\ker D\). Since the remaining
    entries are nonzero scalar multiples of the usual monomials,
    \(E'\) is a basis of \(\mathcal P_3(\mathbb{R})\). Moreover,
    \[
      D(E')=[\,0\ \ 1\ \ x\ \ x^2\,]
        =[\,1\ \ x\ \ x^2\,]
          \begin{pmatrix}
            0&1&0&0\\
            0&0&1&0\\
            0&0&0&1
          \end{pmatrix}.
    \]
    Therefore the matrix of \(D\) with respect to \(E'\) and \(F\) has
    the required form.
\end{enumerate}
\end{proof}
